\section{A uniform schedule in the bit model}
\label{sec:uniform-computation}

The table and parameters of Section~\ref{sec:bounded-suffix} depend only on
the finite instance, but their construction need not be efficient.  We now
specify when that construction is performed and account for all computation
charged to an individual request.  The schedule uses the allowed dependence
on the binary length of the request counter.

\paragraph{Canonical preprocessing.}
Store $n,k,d$ in a fixed canonical encoding: binary integers have no leading
zeroes, each rational distance has a positive denominator and relatively
prime numerator and denominator, and matrix entries occur in their fixed
index order.  Use a fixed unambiguous encoding of lengths and delimiters.
Store the actual labeled tuple $s$, initialize a binary request counter to
zero, and initialize a flag indicating that suffix service has not started.
This takes a uniform polynomial number of bit operations in $L$ and produces
$O(L)$ bits.  For completeness, rational reduction is polynomial: long
division has polynomial bit cost, and in the Euclidean algorithm the larger
positive magnitude decreases by at least a factor of two within two
remainder steps.  Thus each reduction takes polynomially many bit
operations.  The explicit matrix encoding also gives $n\le O(L)$ and hence
$k\le O(L)$.  No table construction is part of preprocessing.

\paragraph{One finite constructor.}
Fix once and for all a deterministic, finite-alphabet bit machine $P$ with
the following behavior on the canonical encoding of $(n,k,d)$.  It computes
$D,\delta,R,F,M,H$, constructs the rational optimum and the dyadic table in
Proposition~\ref{prop:finite-table}, and writes its complete output on an
initially blank output tape.  It outputs the parameters, including $a$ and
$b$, and an explicit list of table records.  Each record contains its key
$(w,r,c)$ and the dyadic probability numerators in label order, with zero
numerators for disallowed labels.  Its strings and records are delimited
unambiguously.  Include, in addition, four unary strings of lengths
$H,R,M,b$.  If $S$ is the total output length in bits, then
\begin{equation}
 H,R,M,b\le S.
 \label{eq:unary-size}
\end{equation}
The output is a consecutively written finite string, so its representation
cannot describe a large blank or compressed table as though that table had
already been written.

All choices in $P$ are deterministic: use length and lexicographic orders
for finite indices and the fixed elimination, back-substitution, and
rounding rules from Section~\ref{sec:finite-tables}.  Exact integer and
rational operations are implemented by ordinary bit algorithms, rather
than counted as single machine steps.  Positivity of $\delta$ and
Proposition~\ref{prop:finite-table} show that $P$ halts on every valid
input.  Denote its finite number of elementary transitions by $T_P$, adding
an initial transition if necessary so that $T_P\ge1$.  Neither $C_0$, a
source policy, nor a bound on $T_P$ is an input to this machine.  In
particular, the existence theorem certifies feasibility of its finite
calculation; it supplies no computational oracle.

\begin{lemma}[Bounded implementation of the suffix]
\label{lem:bounded-implementation}
Given the explicitly written constructor output of length $S$, the suffix
routine can be implemented with a uniform polynomial number of bits of
storage and a uniform polynomial number of bit operations per request in
$L+S+2$.  These bounds hold for every reachable state and every outcome of
its random choices.  Each table sample uses exactly $b$ fresh unbiased
random bits.
\end{lemma}

\begin{proof}
Write $m=L+S+2$ for this proof.  By Equation~\eqref{eq:unary-size}, the numerical
values $H,R,M,b$ are at most $m$, and $n,k\le O(m)$.  Use explicit lists,
binary counters, and bounded strings for the following state.  The actual
and virtual tuples contain $k$ point indices each.  The kept word has at
most $H$ indices.  The filter has at most $H$ live entries by
Lemma~\ref{lem:filter}; each entry contains a tuple of $k$ indices and a
counter between zero and $M$.  Entries with the same data may remain
separate, as required by the filter construction.

For each raw request, scan every live tuple to determine whether it covers
that request.  If the request is kept, copy the surviving entries into a
new list and generate at most $k$ children for each removed entry.  Even an
implementation that temporarily materializes all $kH$ possible children
uses polynomial space: allowing $O(m)$ bits per index or counter gives
$O(m^3)$ bits for the old live list and $O(m^4)$ bits for all temporary
children.  Appending to the kept word and comparing or copying the lists
requires only polynomially many operations on those strings.  The final
list again has at most $H$ entries.  Continue these deterministic filter
updates in fallback mode as well, so the stored kept word is the whole
filtered prefix specified in Section~\ref{sec:bounded-suffix}.

The remaining state is also bounded.  Store the current raw block's set
of at most $k+1$ points, a counter of completed raw blocks bounded by
$R$, and the fallback marks together with their assigned labels.  Mark
sets have at most $k$ protected points; the possible next new point is
handled at that block's endpoint.  A mode flag and the virtual positive
move counter, capped at $M+1$, suffice for switching.  Set membership,
choice of the least eligible label, counter updates, and changes of tuple
entries all use finite scans and copies.  Epoch and block resets erase or
reuse precisely these bounded lists.  No raw epoch length, global epoch
number, cumulative movement total, or unfiltered request history is
stored.

At a table decision, locate the record by scanning the entire output and
comparing its key with the kept prefix, current request, and current
virtual tuple.  The key uses the prefix before appending the current kept
request.  Lemma~\ref{lem:filter} and the table-mode invariant ensure that
every required key has prefix length less than $H$.  All rows, including
rows whose original rational marginal was zero, were written by $P$.
Consequently lookup and the same storage bounds apply also to states
reached only through dyadic rounding.

If the row numerators are $N_1,\ldots,N_k$, draw exactly $b$ unbiased bits
to obtain a uniform integer $U\in\{0,\ldots,2^b-1\}$.  Choose the least
$j$ with $\sum_{i\le j}N_i>U$.  Such a $j$ exists because the numerators
sum to $2^b$; it is allowed because disallowed labels have numerator zero.
Only $k$ additions and comparisons of integers of at most $b+1$ bits are
needed, in addition to the $b$ random bits.  This sampling has no rejection
loop.  On a dropped table-mode request, or in fallback mode, the prescribed
label is found by the bounded scans already described.

These operations have a fixed number of nested loops, each over strings
or lists with polynomially bounded lengths.  Elementary arithmetic uses
bounded binary strings.  Implementing every access, scan, and copy on
ordinary tapes therefore gives a fixed-degree polynomial time bound, with
absolute constants. Use fixed-origin tape regions for the stored lists
and return their heads to those origins during cleanup.
At the end of each request discard or erase obsolete
lists and temporary strings, carrying only the stated compact state into
the next request.  Cleanup is included in the same polynomial bound.
Every bound is deterministic and independent of the probability of the
current state.
\end{proof}

\paragraph{The online schedule.}
At request number $t\ge1$, set
\begin{equation}
 g(t)=\lfloor\log_2(t+1)\rfloor.
 \label{eq:schedule-budget}
\end{equation}
Simulate $P$ from scratch on the stored canonical input for at most $g(t)$
elementary transitions.  If it has not halted, serve the current request
with label $1$.  If it has halted, use its output to serve the current
request by the suffix routine.  At the first successful simulation,
initialize a fresh suffix epoch, including this current request, and set
the activation flag.  At every later success, retain the suffix state
from the preceding request and perform its next step.  Recomputing the
identical table does not reset that state; only the prescribed epoch
boundaries do so.  Erase the simulation workspace and table after use,
preserving the canonical input, actual tuple, counter, activation flag,
and compact suffix state.

Figure~\ref{fig:activation} separates the repeated constructor simulation
from the persistent online state. The output-size estimate is the link
between the simulation budget and the suffix implementation bound.

\begin{proposition}[Uniform scheduling bound]
\label{prop:schedule}
There is one polynomial $p$, independent of the instance, such that
preprocessing uses at most $p(L)$ bit operations and serving request $t$
uses at most
\[
 p\bigl(L+\lceil\log_2(t+1)\rceil\bigr)
\]
bit operations, for every reachable history and internal state.  Activation
occurs at a finite deterministic time $t_*$ depending only on $(n,k,d)$.
\end{proposition}

\begin{proof}
The request counter has $O(g(t)+1)$ bits.  Incrementing it and determining
$g(t)$ costs a polynomial in its length.  A simulation lasting at most
$g(t)$ elementary transitions uses the copied input and at most
$O(g(t)+1)$ additional cells on each of a fixed number of tapes.  Explicit
tape simulation, including copying the $O(L)$-bit input, locating tape
heads, counting transitions, and erasing the workspace, has uniform
polynomial cost in $L+g(t)+2$.  This assertion concerns the actual capped
simulation; it does not charge an uncapped arithmetic subroutine as one
transition.

Whenever a simulation succeeds, a fixed finite-alphabet machine has
written at most a constant times $g(t)+1$ output bits.  Thus
\begin{equation}
 S=O(g(t)+1)
 \label{eq:output-budget}
\end{equation}
before any suffix operation is attempted.  Lemma~\ref{lem:bounded-implementation}
then bounds all suffix work by a uniform polynomial in $L+g(t)+2$.
The failure action, updating the actual tuple, and issuing the chosen
label have smaller polynomial bounds.  The preceding request's budget
was no larger than the current one.  Its compact state and any storage
to be accessed, copied, or erased are consequently bounded by the same
current polynomial.  Our cleanup convention leaves no accumulated
workspace depending on the number of earlier requests.  This also
accounts for the first initialization and every later reset.

Choose absolute constants $A$ and an integer $c$ large enough to dominate
these finitely many polynomial bounds, and set $p(x)=A(x+2)^c$.  Increase
$A,c$ if necessary to cover preprocessing.  Since
$g(t)\le\lceil\log_2(t+1)\rceil$, this is the single polynomial asserted.
The bound is worst case, so it implies the required conditional expected
bound at every reachable state.

Finally, define for analysis only
\begin{equation}
 t_* = \min\{t\ge1:g(t)\ge T_P\}=2^{T_P}-1.
 \label{eq:activation-time}
\end{equation}
It is finite.  The constructor and its canonical input depend only on
$(n,k,d)$, so the same is true of $T_P$ and $t_*$.  In particular,
irrelevant padding or unreduced representations of the original rational
input cannot change this threshold, nor can $s$ or the request sequence.
Monotonicity of $g$ ensures that every simulation from $t_*$ onward
succeeds.  The algorithm never computes $T_P$ or $t_*$ in advance; it
detects success by the ordinary halting state of its capped simulation.
\end{proof}

\begin{figure}[tbp]
\centering
\begin{tikzpicture}[
  font=\small,
  box/.style={draw=black!65,rounded corners=2pt,align=center,
    inner sep=6pt,text width=3.1cm,minimum height=1.05cm},
  flow/.style={-{Latex[length=2mm]},thick}
]
\node[box,fill=black!3] (input) at (0,0)
  {Canonical $(n,k,d)$\\Request counter $t$};
\node[box,fill=blue!5] (simulate) at (4.7,0)
  {Capped simulation\\of constructor $P$\\$g(t)$ transitions};
\node[box,fill=orange!8] (failure) at (9.4,0)
  {Serve with label $1$};
\node[box,fill=blue!5] (output) at (4.7,-2.4)
  {Completed table\\and parameters\\$S=O(g(t)+1)$};
\node[box,fill=blue!5] (suffix) at (9.4,-2.4)
  {First success:\\activate\\Later successes:\\advance state};
\node[align=center,text width=3.7cm,font=\footnotesize] (state) at (9.4,-4.05)
  {Retain suffix state\\for the next request.};
\node[align=left,text width=3.6cm] at (0,-2.4)
  {$g(t)=\lfloor\log_2(t+1)\rfloor$\smallskip\\
   First success:\\$t_*=2^{T_P}-1$.};
\draw[flow] (input) -- (simulate);
\draw[flow] (simulate) -- node[above,align=center,font=\footnotesize]{not\\halted} (failure);
\draw[flow] (simulate) -- node[right,font=\footnotesize]{halted} (output);
\draw[flow] (output) -- (suffix);
\draw[flow] (suffix.south) -- (state.north);
\end{tikzpicture}
\caption{Each request pays for a fresh capped simulation. Before activation,
only one fixed label serves requests. After activation, the complete table's
explicit size fits the current budget, and the suffix routine advances its
persistent state, with resets only at its prescribed epoch boundaries.
The construction time $T_P$ and activation time $t_*$ are
finite instance-dependent quantities; their magnitude enters the additive
movement term.}
\label{fig:activation}
\end{figure}


\begin{proof}[Proof of Theorem~\ref{thm:main}]
The preprocessing and online schedule above define a uniform causal
algorithm using only unbiased random bits.  Every decision serves the
current request by moving one labeled actual server.  By
Proposition~\ref{prop:schedule}, its bit complexity has the required bound
with one polynomial $p$.

Fix any finite request sequence $\sigma$ independent of the random bits.
The requests before activation cost at most
$(t_*-1)D$ in total, since each moves a single server over distance at most
$D$.  If $\sigma$ ends before activation, this already bounds its entire
cost.  Otherwise Proposition~\ref{prop:suffix}, applied to the suffix
starting at the deterministic time $t_*$ and compared with the optimum
from $s$ for the whole sequence, gives
\begin{equation}
 \E[\operatorname{cost}_A(\sigma)]
 \le (4a+2)\operatorname{OPT}_s(\sigma)
       +(t_*-1)D+(2a+2)kD.
 \label{eq:full-movement-bound}
\end{equation}
The same inequality holds for shorter sequences, including the empty
sequence, by nonnegativity of its right-hand terms.

Since $a\le C_0\log^2(k+1)$ and $\log(k+1)\ge\log3$, set
\begin{equation}
 C=4C_0+\frac{2}{(\log3)^2},
 \qquad
 B_{d,k,s}=(t_*-1)D+(2a+2)kD.
 \label{eq:final-constants}
\end{equation}
Then $4a+2\le C\log^2(k+1)$, proving the desired movement inequality.
The constant $C$ is absolute, and $B_{d,k,s}$ is finite and depends only
on the instance.  Its formula happens not to require $s$.  Repeated
positions in $s$ cause no difficulty: the suffix comparison allowed
arbitrary actual starting tuples, while the existence theorem and finite
table were used only from the distinct virtual tuple $u$.
\end{proof}
